\chapter{Introduction}
\label{ch:intro}

\section{The problem}

An office accumulates idle computation. Desktops sit at a few percent
utilisation for most of the working day and at zero overnight; mini-PCs bought
for one task spend the rest of their lives idle. The aggregate is substantial
--- this pool alone has eight four-core mini-PCs on the network, with more
machines joining --- and none of it is reachable by anyone who needs compute.

Batch systems solve exactly this, and HTCondor was built for precisely this case:
harvesting cycles from machines whose owners have first claim on them. The
engine is not the problem.

The problem is the interaction model that conventionally comes with it.

\begin{keypoint}
\textbf{The requirement that shaped every decision here:} users must never have
to log into a submission host. Not because logging in is difficult, but because
the moment it is required, the batch system becomes a separate place people have
to visit --- with its own home directory, its own environment, its own files to
copy back and forth --- rather than an extension of the machine they already
work on.
\end{keypoint}

That is the friction that keeps academic clusters underused: not the scheduler,
but the ritual around it. A colleague with a script that would finish in ten
minutes across the pool will run it for six hours locally rather than learn
another machine's filesystem layout.

\section{Design goals}

\begin{table}[H]
\centering
\begin{tabularx}{\textwidth}{L{0.42\textwidth} Y}
\toprule
\textbf{Goal} & \textbf{Consequence} \\
\midrule
Users never log into the entry node & Remote submit only; no shells, no home directories on the entry node \\
Users' own \emph{desktops} are also workers & The pool is a symmetric co-operative, not a donated cluster. Laptops may be submit-only \\
The owner always wins on their own machine & Preemption policy and machine \kn{RANK} are mandatory, not optional \\
Jupyter is allowed, but never on the entry node & Notebooks spawn as Condor jobs on workers \\
Mixed hardware is the normal case & Two architectures, three distributions, some GPUs \\
\bottomrule
\end{tabularx}
\caption{Design goals and what each one forces.}
\end{table}

The second goal is the one that makes this design unusual. In a conventional
pool, workers are dedicated machines and submitters are separate. Here every
participant contributes their own desktop and submits from it. That symmetry is
the adoption mechanism --- \emph{I contribute and I still get my PC first} ---
and it is also the source of most of the engineering difficulty in this report,
because a worker that somebody is actively using cannot be treated like a
compute node.

\section{Why not the obvious alternatives}

\textbf{Slurm} assumes a consistent user namespace across the cluster. On this
fleet the same login has a different UID on different hosts, which is irrelevant
to HTCondor and fatal to Slurm.

\textbf{Kubernetes} solves container orchestration for services, not opportunistic
batch on machines whose owners reclaim them without warning. The eviction model
is wrong: pods are rescheduled on capacity pressure, not on the arrival of a
human at a keyboard.

\textbf{A dedicated cluster} is the answer if the hardware exists. It does not.
The premise is that the machines are already bought, already deployed and already
doing something else most of the time.

HTCondor's ClassAd matchmaking, its explicit owner policy, and its indifference
to UID consistency make it the correct engine. What follows is almost entirely
about the layer above it: the policy that decides when a machine will accept
work, and the naming that lets machines find each other while moving between
networks.

\section{What was built}

A three-machine pool spanning two subnets with no route between them:

\begin{itemize}
  \item An \textbf{entry node} --- a four-core mini-PC --- carrying the collector,
        negotiator, schedd and a limited startd.
  \item A \textbf{flagship worker} --- sixteen threads, 31\,GB --- the pool's
        only high-core machine.
  \item An \textbf{in-use desktop} --- a four-core mini-PC belonging to someone
        who uses it all day, on the other subnet.
  \item A \textbf{user PC} --- a four-core mini-PC on both subnets that
        \emph{submits and executes}. It is the first machine of the class the
        whole design targets, and the first whose owner has a Condor identity,
        so \kn{RANK} is finally in use.
\end{itemize}

Jobs submitted from the user PC distribute across all four machines and return
their output. Workers carry no network configuration at all: each derives at
startup which of the entry node's two addresses it can reach, so a machine that
moves between subnets reconfigures itself.

\begin{keypoint}
The user PC is the point of the exercise, not an extra worker. It runs the
submit client and a startd and \textbf{no schedd}: its owner submits against
the entry node's queue from their own desktop, and the same desktop executes
other people's jobs when they are not using it. That is the symmetric
co-operative in Table~1.1, and everything before it was scaffolding.

Donation is expected of \emph{desktops}. A laptop that only submits is a
normal member of the pool, not an opt-out --- see \S\ref{sec:notallworkers}.
\end{keypoint}
